\chapter{Точная LCF-проверка усреднение по калибровочной группе Краус-моста}
\label{ch:v8-gauge-twirl-kraus-lcf-migration-gate}

\section{Полный типизированный мультиплет}

Исходный гейт проверял калибровочную ковариантность двенадцатью случайными
преобразованиями. В eDSL физические семейства
$Q_L\leftrightarrow Y_R$ и $X_L\leftrightarrow d_R$ восстановлены как
точные матрицы $\mathbb C^{11}\to\mathbb C^{10}$. После вещественного
разложения они дают соответственно $6+6$ самосопряжённых операторов
\begin{equation}
 D_e=\begin{pmatrix}0&e^*\\e&0\end{pmatrix}
 \in\operatorname{End}(\mathbb C^{21}).
\end{equation}
Их коэффициенты состоят только из целых чисел, $i$ и точной нормировки
$1/\sqrt2$.

Конструктор GKSL принимает нулевой гамильтониан и единичные скорости:
\begin{equation}
 \mathcal L_{q\ell}(Z)
 =\sum_{a=1}^{12}\left(D_aZD_a-\frac12\{D_a^2,Z\}\right).
 \label{eq:v8-lcf-gauge-twirl-generator}
\end{equation}
Точно получено $\mathcal L_{q\ell}(I_{21})=0$.

\section{Калибровочная замкнутость и независимость базиса}

Ядро проверило все восемь эрмитовых генераторов $su(3)$, три генератора
$su(2)$ и точный рациональный гиперзаряд. Для каждого $G$ и каждого $D_a$
коммутатор $i[G,D_a]$ точно восстанавливается в том же двенадцатимерном
вещественном базисе. Матрица этого действия вещественна и кососимметрична,
а общий остаток замыкания равен нулю.

Совместное ядро двенадцати матриц действия имеет размерность
\begin{equation}
 \dim\bigcap_G\ker R_G=0.
 \label{eq:v8-lcf-no-linear-gauge-singlet}
\end{equation}
Поэтому линейного калибровочного синглета нет точно, а не только после выборочного
усреднения. Одновременно равенство скоростей позволяет применить тождество
$\sum_bO_{ba}O_{bc}=\delta_{ac}$: квадратичная сумма
\eqref{eq:v8-lcf-gauge-twirl-generator} неизменна при любом вещественном
ортогональном изменении Краус-базиса.

\section{Точное снятие двухсекторности}

На базисе $P_q/\sqrt{12},P_\ell/3$ каждая цветная семья даёт
\begin{equation}
 M_0=\begin{pmatrix}
 \frac12&-\frac1{\sqrt3}\\
 -\frac1{\sqrt3}&\frac23
 \end{pmatrix},
 \qquad \Spec M_0=\left\{0,\frac76\right\}.
\end{equation}
Их сумма равна
\begin{equation}
 M_{q\ell}=\begin{pmatrix}
 1&-\frac2{\sqrt3}\\
 -\frac2{\sqrt3}&\frac43
 \end{pmatrix},
 \qquad
 \det(\lambda I-M_{q\ell})=\lambda\left(\lambda-\frac73\right).
 \label{eq:v8-lcf-cross-central-matrix}
\end{equation}
Следовательно, прежняя неподвижная $\mathbb C^2$ сокращается ровно до
одной скалярной линии. Контрольные внутренние лептонные семейства дают
нулевую матрицу $2\times2$.

Для произвольных $\gamma_1,\gamma_2>0$ ограничение равно
$(\gamma_1+\gamma_2)M_0$, поэтому одномерность ядра доказана символически
для всех положительных скоростей. Прежний случайный перебор из 64 точек больше
не является частью доказательства.

\begin{theorem}[Точный квадратичный межсекторный мост]
Полный двенадцатимерный вещественный Краус-мультиплет замкнут под
$su(3)\oplus su(2)\oplus u(1)$, не содержит линейного инвариантного
направления, но его базисно-независимый GKSL-генератор сокращает
$\mathbb CP_q\oplus\mathbb CP_\ell$ до $\mathbb CI$. Результат сохраняется
для любых двух положительных скоростей семейств.
\end{theorem}

\begin{nogocriterion}[Граница миграции]
LCF-сертификат не выводит абсолютную или относительную скорость процесса и
не доказывает происхождение генератора из общего родительского действия. Нельзя
считать физическое взаимодействие полностью выведенным до точной миграции
родительского функционала и его гессиана.
\end{nogocriterion}

\begin{protocol}
\begin{enumerate}
 \item размерность носителя концевых секторов: \textbf{$21$};
 \item вещественных межсекторных скачки: \textbf{$12$};
 \item точных калибровочных генераторов: \textbf{$12$};
 \item остаток замыкания репера скачков: \textbf{ноль};
 \item линейная инвариантная размерность: \textbf{$0$};
 \item независимость от ортогонального Краус-базиса: \textbf{точно};
 \item центральный ненулевой темп: \textbf{$7/3$};
 \item центральное ядро: \textbf{$1$};
 \item положительные скорости: \textbf{символически все};
 \item родительского действия и его гессиан: \textbf{ещё открыты}.
\end{enumerate}
\end{protocol}

Машинный сертификат сохранён в
\path{s2t_v8_gauge_twirl_kraus_lcf_migration_gate_results.json}.